%% Bac de rétention en tôle pliée (exercice guidé de la compétence « Dérivation »).
%% À gauche : la plaque carrée de 30 cm de côté, les quatre carrés de coin de côté x hachurés
%% (ils sont découpés) et les quatre plis en tirets. À droite : le bac obtenu, sans couvercle,
%% en perspective cavalière — hauteur x ; le côté du fond (30 - 2x) est dans \rappel :
%% c'est ce que l'exercice guidé demande de trouver, la variante muette le laisse en blanc.
%% Aucune expression du volume n'est écrite : c'est la réponse attendue de l'élève.
\documentclass[tikz,border=4pt]{standalone}
\usepackage{liaisons}
\begin{document}
\begin{tikzpicture}[x=1cm,y=1cm,
  cote/.style={{Stealth[length=4pt]}-{Stealth[length=4pt]}, line width=0.6pt, black!75},
  attache/.style={line width=0.4pt, black!55},
  pli/.style={line width=0.8pt, dash pattern=on 3pt off 2pt, solideB},
  etiq/.style={font=\small},
  petit/.style={font=\scriptsize, text=black!65},
  titre/.style={font=\scriptsize\bfseries, align=center}]

\def\S{4.2}     % côté de la plaque
\def\c{0.85}    % côté du carré de coin

%% ==================== (a) la plaque à plat ==============================
%% fond du bac (carré central) et bandes qui seront relevées
\fill[solideE!12] (\c,\c) rectangle (\S-\c,\S-\c);
\fill[solideD!12] (\c,0) rectangle (\S-\c,\c);
\fill[solideD!12] (\c,\S-\c) rectangle (\S-\c,\S);
\fill[solideD!12] (0,\c) rectangle (\c,\S-\c);
\fill[solideD!12] (\S-\c,\c) rectangle (\S,\S-\c);

%% les quatre carrés de coin : découpés
\hachures[draw=solideA!70]{(0,0) rectangle (\c,\c)}
\hachures[draw=solideA!70]{(\S-\c,0) rectangle (\S,\c)}
\hachures[draw=solideA!70]{(0,\S-\c) rectangle (\c,\S)}
\hachures[draw=solideA!70]{(\S-\c,\S-\c) rectangle (\S,\S)}

%% contour de la plaque et lignes de pli
\draw[line width=0.9pt, line join=round] (0,0) rectangle (\S,\S);
\draw[pli] (\c,0) -- (\c,\S)  (\S-\c,0) -- (\S-\c,\S)
           (0,\c) -- (\S,\c)  (0,\S-\c) -- (\S,\S-\c);

%% cote du côté de la plaque
\draw[attache] (0,0) -- (0,-0.72);
\draw[attache] (\S,0) -- (\S,-0.72);
\draw[cote] (0,-0.55) -- (\S,-0.55);
\node[etiq, fill=white, inner sep=1.5pt] at (\S/2,-0.55) {$30$ cm};

%% cote du carré de coin (en haut à gauche)
\draw[attache] (0,\S) -- (0,\S+0.62);
\draw[attache] (\c,\S) -- (\c,\S+0.62);
\draw[cote] (0,\S+0.45) -- (\c,\S+0.45);
\node[etiq, text=solideA, anchor=south, inner sep=2pt] at (\c/2,\S+0.52) {$x$};

%% cote du même carré, en hauteur (en bas à gauche)
\draw[attache] (0,0) -- (-0.62,0);
\draw[attache] (0,\c) -- (-0.62,\c);
\draw[cote] (-0.45,0) -- (-0.45,\c);
\node[etiq, text=solideA, anchor=east, inner sep=2pt] at (-0.52,\c/2) {$x$};

\node[titre, anchor=north] at (\S/2,-1.05) {(a) la plaque découpée};
\node[petit, anchor=north, align=center] at (\S/2,-1.42)
  {carrés hachurés : découpés\\ tirets : lignes de pli};

%% ==================== (b) le bac obtenu =================================
\begin{scope}[shift={(6.6,0.75)}]
  \def\b{3.0}       % côté du fond
  \def\h{1.0}       % hauteur (= x)
  \def\dx{0.85}     % fuyante
  \def\dy{0.55}

  \coordinate (F1) at (0,0);
  \coordinate (F2) at (\b,0);
  \coordinate (F3) at (\b+\dx,\dy);
  \coordinate (F4) at (\dx,\dy);
  \coordinate (T1) at (0,\h);
  \coordinate (T2) at (\b,\h);
  \coordinate (T3) at (\b+\dx,\h+\dy);
  \coordinate (T4) at (\dx,\h+\dy);

  %% intérieur : le fond, puis les deux parois du fond de la vue
  \filldraw[fill=solideE!18, draw=black!65, line width=0.7pt, line join=round]
     (F1) -- (F2) -- (F3) -- (F4) -- cycle;
  \filldraw[fill=solideD!10, draw=black!65, line width=0.7pt, line join=round]
     (F1) -- (F4) -- (T4) -- (T1) -- cycle;
  \filldraw[fill=solideD!10, draw=black!65, line width=0.7pt, line join=round]
     (F4) -- (F3) -- (T3) -- (T4) -- cycle;
  %% extérieur : paroi avant, puis paroi droite
  \filldraw[fill=solideD!25, draw=black!75, line width=0.8pt, line join=round]
     (F1) -- (F2) -- (T2) -- (T1) -- cycle;
  \filldraw[fill=solideD!20, draw=black!75, line width=0.8pt, line join=round]
     (F2) -- (F3) -- (T3) -- (T2) -- cycle;
  %% bord supérieur (le bac est ouvert)
  \draw[line width=0.8pt, line join=round] (T1) -- (T2) -- (T3) -- (T4) -- cycle;

  %% cote de la hauteur
  \draw[attache] (T2) -- ++(0.95,0);
  \draw[attache] (F2) -- ++(0.95,0);
  \draw[cote] (\b+0.78,0) -- (\b+0.78,\h);
  \node[etiq, text=solideA, anchor=west, inner sep=2pt] at (\b+0.84,\h/2) {$x$};

  %% cote du côté du fond
  \draw[attache] (F1) -- (0,-0.72);
  \draw[attache] (F2) -- (\b,-0.72);
  \draw[cote] (0,-0.55) -- (\b,-0.55);
  \node[etiq, fill=white, inner sep=1.5pt] at (\b/2,-0.55) {\rappel{$30-2x$}};

  %% les deux légendes sont alignées sur celles de la plaque (décalage de la scène : 0,75)
  \node[titre, anchor=north] at (\b/2,-1.80) {(b) le bac, sans couvercle};
  \node[petit, anchor=north, align=center] at (\b/2,-2.17)
    {le fond est un carré,\\ les quatre bandes sont relevées};
\end{scope}
\end{tikzpicture}
\end{document}
